\section{LLM-as-a-Judge and Inter-Judge Agreement}
\label{sec:sota-judge}

The methodology this thesis borrows was developed outside software engineering.
G-Eval established criterion-based scoring by a language
model~\cite{liu2023geval}; MT-Bench extended it to whole responses and measured the
judge against human preference~\cite{zheng2023judging}; Chatbot Arena scaled
pairwise human comparison with model identities hidden to a size at which
rankings stabilise~\cite{chiang2024chatbotarena}. That literature also produced the
countermeasures now considered good practice: several judges rather than one,
coverage across model families so that no single provider's idiosyncrasy dominates,
blinding and pairwise presentation to control position and identity effects, and
explicit reporting of inter-judge agreement so that a reader can see how reliable
the instrument was. Panickssery et al.\ supplied the sharpest reason for the
cross-family requirement by demonstrating that judges recognise and favour their
own generations~\cite{panickssery2024selfpreference}.

He et al.\ survey the migration of these methods into software engineering and find
adoption to be uneven~\cite{he2025llmasjudgese}. The code-review evaluations closest
to this thesis illustrate the point: DeepCRCEval scores its rubric with a single
judge~\cite{hong2024deepcrceval}, and CRScore is a single metric
pipeline~\cite{naik2025crscore}. Neither reports an agreement coefficient between
independent scorers of the same reviews, because neither has two.

This matters more for a context-augmentation study than for a model-ranking study,
for a reason specific to the intervention. Adding retrieved material to a prompt
reliably makes the generated review longer. Verbosity bias, in which a judge
rewards length independent of content, is among the best-documented judge
pathologies~\cite{zheng2023judging}. A single-judge protocol cannot separate a real
improvement from a judge preferring more text, and self-preference compounds the
problem when the generator and the judge come from the same family. Bringing the
panel methodology into code-review evaluation and measuring how much the
conclusion depends on individual judges is therefore
not a methodological flourish here; it is what makes the central claim
interpretable, and Chapter~\ref{chap:results} reports the difference it makes.
